\documentclass[11pt]{article}

%==============================================================================
%  LPS / PS-LPS Program: Plan of Action and Phase Sequencing  (living plan)
%  Reframed as an algebra: two base estimators (LPS, LCov) crossed with
%  orthogonal refinement axes (prediction synchronization; the local
%  dimension/scale ladder), plus a cross-cutting anchoring choice.
%==============================================================================

\usepackage[margin=1in]{geometry}
\usepackage{amsmath,amssymb}
\usepackage{xcolor}
\usepackage{booktabs}
\usepackage{array}
\usepackage{enumitem}
\usepackage{microtype}
\microtypesetup{expansion=false}
\usepackage{caption}
\usepackage[most]{tcolorbox}
\usepackage{datetime2}
\DTMsetup{showzone=false}
\usepackage[hidelinks]{hyperref}
\widowpenalty=10000
\clubpenalty=10000
\raggedbottom

\definecolor{ink}{HTML}{24303D}
\definecolor{accent}{HTML}{2563EB}
\definecolor{accent3}{HTML}{15803D}
\setlist{topsep=3pt,itemsep=2pt,parsep=1pt}

\newtcolorbox{gate}{
  enhanced, colback=accent3!5, colframe=accent3!55!black, boxrule=0.4pt, arc=2pt,
  left=7pt, right=7pt, top=4pt, bottom=4pt, title=Go/no-go gate,
  fonttitle=\bfseries\sffamily\small, coltitle=accent3!45!black,
  attach title to upper={\quad}}
\newtcolorbox{note}{
  enhanced, colback=accent!4, colframe=accent!55!black, boxrule=0.4pt, arc=2pt,
  left=7pt, right=7pt, top=4pt, bottom=4pt, title=Why this ordering,
  fonttitle=\bfseries\sffamily\small, coltitle=accent!50!black,
  attach title to upper={\quad}}

\newcommand{\field}[1]{\smallskip\noindent\textbf{\textsf{#1.}}\ }

\title{\textbf{LPS / PS-LPS Program: Plan of Action and Phase Sequencing}\\[3pt]
\large \textnormal{Living roadmap --- refined as a base $\times$ refinement algebra}\\[4pt]
\normalsize A dependency-ordered roadmap: two base estimators (LPS, local
covariance), the prediction-synchronization and local dimension/scale refinement
axes, the anchoring choice, the simplicial-complex integration, and the occupancy
capstone}
\author{\small Internal program plan --- LPS / \texttt{geosmooth} program\\[2pt]
\footnotesize source: \nolinkurl{~/current_projects/geosmooth/dev/programs/lps_lcov_omics_program/lps_lcov_omics_program_plan.tex}\\[1pt]
\footnotesize stable canonical source; version history is tracked by git}
\date{\DTMnow\ \textnormal{\footnotesize ET}}

\begin{document}
\maketitle

\begin{abstract}\noindent
The first plan listed the program as a linear sequence of phases. Working through
the local-chart geometry made a cleaner structure visible: the program is a small
\emph{algebra}. There are two \textbf{base estimators} --- LPS (the local-polynomial
conditional mean) and LCov (the local conditional covariance/association) --- and a
short list of \textbf{refinement operators} that apply to either base and
\emph{compose}. The operators fall on two orthogonal axes: \textbf{(A) prediction
synchronization} (PS), which couples overlapping local models across anchors, and
\textbf{(B) the local dimension/scale ladder}, which sets each anchor's local chart.
Axis~B has four rungs of increasing structure --- fixed dimension/scale (the base);
\textbf{LDS}, the joint per-anchor estimation of \emph{both} local dimension and
scale, unregularized; \textbf{LDR}, regularization of the dimension field with scale
held fixed (the original framing); and \textbf{LDSR}, co-regularization of dimension
\emph{and} scale together. A third, cross-cutting axis is the \textbf{anchoring}
choice: local models centered at the data points, or at a regular/derived grid. The
two refinement axes are genuinely orthogonal --- and they \emph{meet} at the chart
nerve, which is therefore named as the integration phase rather than a sub-part of
any one operator. The validated base emits the reusable $S$/df/LOO/band toolkit; the
synthetic-data enabler runs throughout; everything converges into the nerve
integration and then the occupancy / quasi-equilibrium capstone and the terminal
16S application. The whole program runs under the two-agent audit-independence
workflow. As of this revision, Tier-0 (the LPS base) is validated and merged, and the
Tier 1--4 estimator gates are in integration.
\end{abstract}

{\small\tableofcontents}

%==============================================================================
\section{How to read this, and the one rule that wraps every phase}\label{sec:read}
%==============================================================================

This is a sequencing document, not a methods document; the methods live in the
design notes referenced at the end, and each phase points to its note. The job here
is to say what is built first, in what combinations, and why the order is forced by
dependencies rather than preference.

One rule wraps every phase: each is executed under the \emph{two-agent
audit-independence workflow} --- an implementer builds; an independent auditor
validates against a fixed charter (the auditor's scope is not set by the
implementer; the handoff carries evidence and admitted limitations but not the
audit's questions; ``report rendered'' is never ``phase accepted''; and a gate is
real only if it can be made to fail). A phase is not complete --- and dependents do
not start --- until its go/no-go gate is passed by that independent audit or
adjudicated by the investigator.

%==============================================================================
\section{The map: two bases, two refinement axes, one anchoring choice}\label{sec:map}
%==============================================================================

\field{Two base estimators} \textbf{LPS} estimates the conditional mean
$\mathbb{E}[y\mid x]$ by local polynomial regression in a local chart. \textbf{LCov}
estimates the local conditional second-order structure (covariance / precision /
association) in the same chart geometry. They share the chart machinery; they differ
in what is fit on the chart.

\field{Axis A --- prediction synchronization (PS)} An operator that couples
\emph{overlapping} local models: where two anchors' charts share data, their
predictions are regularized toward agreement (prediction-overlap coupling, strength
$\lambda_{\mathrm{sync}}$). PS acts \emph{across} anchors, on the overlap graph. It
applies to LPS (the existing PS-LPS) and, by analogy, to LCov.

\field{Axis B --- the local dimension/scale ladder} An operator that sets
\emph{each} anchor's local chart --- its intrinsic dimension $d_i$ and its scale
$K_i$ (support size / bandwidth). Four rungs:
\begin{description}[leftmargin=1.1em,style=nextline]
\item[base] $d$ and $K$ fixed (a numeric \texttt{chart.dim} and one CV-selected
global $K$): the current default.
\item[LDS --- local dimension \emph{and} scale, unregularized] Estimate \emph{both}
fields per anchor, jointly. This rung is itself two coupled directions --- dimension
and scale --- because the scale at which a neighborhood is read determines its
apparent dimension and vice versa. No spatial smoothing: the honest raw estimator.
\item[LDR --- local dimension regularization] Regularize the dimension field
($\ell_1$/TV over a graph) while holding scale fixed/global. This is the
program's \emph{original} way of thinking about local dimension; it is the special
case of the next rung with scale-regularization switched off.
\item[LDSR --- local dimension \emph{and} scale regularization] Co-regularize both
fields together over the chart nerve, subject to the points-per-parameter coupling.
The full object; it nests LDR.
\end{description}

\field{Axis C (cross-cutting) --- anchoring} Where the local models are centered:
at the \textbf{data points} (the current choice), or at a \textbf{regular/derived
grid} that decouples the anchor set from the sample (\S\ref{sec:grid}). This choice
crosses every cell of the base$\times$refinement product.

\field{The product} The named variants are the cells of
$\{\text{LPS},\text{LCov}\}\times\{\text{base},\,\text{PS},\,\text{LDS},\,\text{LDR},
\,\text{LDSR},\,\text{PS+LDSR}\}$, each under a choice of anchoring:

\begin{center}\small
\renewcommand{\arraystretch}{1.25}
\begin{tabular}{@{}l l l@{}}
\toprule
\textbf{Refinement} & \textbf{base $=$ LPS} & \textbf{base $=$ LCov}\\
\midrule
none & LPS & LCov\\
$+$ PS (synchronization) & LPS+PS \ (\,$=$ PS-LPS\,) & LCov+PS\\
$+$ LDS (local dim \& scale, raw) & LPS+LDS & LCov+LDS\\
$+$ LDR (dim regularization, scale fixed) & LPS+LDR & LCov+LDR\\
$+$ LDSR (dim \& scale co-regularization) & LPS+LDSR & LCov+LDSR\\
$+$ PS $+$ LDSR (full) & LPS+PS+LDSR & LCov+PS+LDSR\\
\bottomrule
\end{tabular}
\end{center}

\begin{note}
The two refinement axes are \emph{orthogonal}: PS chooses how anchors talk to each
other; the dimension/scale ladder chooses what each anchor's chart is. That is why
they compose into LPS+PS+LDSR rather than competing --- and why both ultimately live
on the same object, the nerve of the chart cover (\S\ref{sec:nerve}): PS's overlap
graph and LDSR's co-regularization graph are the \emph{same} graph. The nerve is
therefore the integration phase, not a sub-part of either operator.
\end{note}

%==============================================================================
\section{Phase 0 --- the LPS base (the foundation that emits the tools)}\label{sec:phase0}
%==============================================================================

\field{Objective} Prove the base smoother is what it claims, and produce the
reusable estimator toolkit every refinement reuses.
\field{Inputs} None (the foundation).
\field{Deliverables} The Tier-0 correctness battery (polynomial reproduction; the
linear-smoother identity $\hat y = Sy$ with $\mathrm{tr}\,S$ and analytic LOO/GCV;
boundary behaviour; consistency; binary recovery and calibration; CV no-leakage;
degenerate geometry), then the Tier 1--4 plan (honest nested/grouped CV; bandwidth
and kernel semantics; binary-mode hygiene; uncertainty bands). Spec: the experimental
plan note.
\field{Gate} \begin{gate} The Tier-0 battery passes deterministically and the
linear-smoother toolkit ($S$, df, LOO/GCV, bands) is implemented and tested. This is
both the correctness gate and the \emph{source} of the tools the refinements
require: df-matched comparisons (PS), per-anchor local GCV (the dimension/scale
ladder), and bands (LCov, capstone).\end{gate}
\field{Status (this revision)} \emph{Met.} Tier-0 is validated and merged to
\texttt{main}; the $S$/df/LOO/band toolkit exists (\texttt{R/lps\_uncertainty.R}).
Tier 1--4 gates (bandwidth multiplier; nested/grouped CV; binary hygiene and the
ridge-alignment and NA-consistency fixes; pointwise bands) are in integration.

%==============================================================================
\section{Axis B --- the local dimension/scale ladder}\label{sec:ladder}
%==============================================================================

This axis is the heart of the refinement. It is developed bottom-up: build the raw
joint estimator (LDS) first, then the two regularizations (LDR, LDSR). Spec: the
implementation notes (dimension) and the per-point-scale sketch.

\subsection{LDS --- joint local dimension and scale (the estimator)}\label{sec:lds}
\field{Objective} Estimate, per anchor, \emph{both} the intrinsic dimension $d_i$
and the scale $K_i$ (support / bandwidth), recognizing that they are coupled: the
locally MSE-optimal scale depends on local curvature, density, noise, and $d_i$, and
$d_i$ is read off the neighborhood spectrum at scale $K_i$.
\field{Inputs} Phase 0 (the chart machinery and the $S_{ii}$/LOO toolkit); the toy
synthetic generators.
\field{Deliverables} The joint per-anchor $(d_i,K_i)$ fields. \emph{One design fork
decides whether the Phase-0 toolkit survives unchanged:} an \textbf{X-only} scale
(geometry/density-driven; keeps $\hat y = Sy$ linear and the df/LOO/band toolkit
exact) versus a \textbf{y-aware} scale (local GCV/LOO or Lepski; better MSE but
$S$ becomes $y$-dependent, so df/bands move to the selection-aware path). The honest
caveat: response-removal LOO $\neq$ training-removal LOO for a $k$-NN bandwidth, so
the cheap $S_{ii}$ shortcut is a biased proxy for \emph{scale} selection (favour
Lepski, or a leave-a-block-out, for the scale field).
\field{Gate} \begin{gate} On a known varying-scale / varying-dimension DGP, the raw
$(d_i,K_i)$ fields track the truth (smaller $K$ in finer regions; $d_i$ matching the
stratum dimension), and the X-only variant preserves the $\hat y = Sy$ identity
exactly. The fields will be high-variance --- that is expected and is what the next
two rungs address.\end{gate}
\field{Unblocks} LDR, LDSR.

\subsection{LDR --- local dimension regularization (scale fixed)}\label{sec:ldr}
\field{Objective} Stabilize the per-anchor dimension field by $\ell_1$/total-variation
regularization over a graph, holding scale at the global value. (The program's
original local-dimension picture.)
\field{Inputs} Phase 0; LDS's raw dimension field; the toy generators.
\field{Deliverables} The exact graph-cut/TV solve; the test battery T1--T9
(graph-cut exactness, $\ell_1$-keeps-the-boundary vs $\ell_2$-blurs, denoising, the
boundary safeguard, scale/persistence, the coordinate trap, the support cap and soft
face, the frontier, downstream Truth-RMSE). Spec: the implementation notes.
\field{Gate} \begin{gate} The implementation notes' go/no-go: \emph{(the nominal$-$
effective dimension gap is large)} $\wedge$ \emph{(the field improves Truth-RMSE on
the stratum-mixture generator)}, without blurring a true dimension boundary. If
either fails, the within-face field is not promoted; the cheaper soft-face-and-cap
refinement is kept.\end{gate}
\field{Unblocks} The nerve integration; it is also the scale-frozen baseline LDSR
must beat.

\subsection{LDSR --- co-regularized dimension and scale}\label{sec:ldsr}
\field{Objective} Co-regularize the $(d,K)$ fields jointly over the chart nerve,
subject to the points-per-parameter coupling, so scale adapts where the data's
natural scale changes and the two fields stay spatially coherent.
\field{Inputs} LDS (the raw fields); LDR (the dimension-field solver and its
$\ell_1$-boundary property); the chart nerve (\S\ref{sec:nerve}).
\field{Deliverables} The coupled objective
$\sum_i \ell_i(K_i,d_i) + \lambda_K\,\mathrm{TV}_G(K) + \lambda_d\,\mathrm{TV}_G(d)$
subject to $\binom{d_i+p}{p}\le K_i$, solved on the nerve graph $G$ with the LDR
machinery; identifiability guards (the points-per-parameter constraint; a parsimony
tie-break; never selecting scale on in-sample RSS). Spec: the per-point-scale sketch.
\field{Gate} \begin{gate} LDSR beats LDR (scale-frozen) at \emph{matched df} on
heteroscale DGPs, and reduces to $\approx$ the global $K$ on a homogeneous truth (no
spurious variation); the $\ell_1$ co-regularization preserves a true joint
scale/dimension boundary that $\ell_2$ would smooth away.\end{gate}
\field{Unblocks} LPS+PS+LDSR; the LCov ladder; the capstone's basin geometry.

%==============================================================================
\section{Axis A --- prediction synchronization (PS / PS-LPS)}\label{sec:ps}
%==============================================================================

\field{Objective} Validate, and put on firm footing, the synchronization of
overlapping local models through prediction-overlap regularization.
\field{Inputs} Phase 0 (the validated base \emph{and} the toolkit). PS is independent
of the dimension/scale ladder --- it can be developed in parallel.
\field{Deliverables} Exploit PS-LPS's linearity in the response: exact effective
df $=\mathrm{tr}\,S(\lambda_{\mathrm{sync}})$ and analytic LOO/GCV in place of
$K$-fold CV; compare to LPS at \emph{matched df} (not matched $\lambda$); fix the
ridge-path discontinuity at $\lambda_{\mathrm{sync}}=0$; make the synchronization
graph and weights explicit; characterize the estimand and the
$\lambda_{\mathrm{sync}}\to\infty$ limit; add bands. Spec: the audit memo's PS-LPS
sections and the chart-aware synchronization.
\field{Gate} \begin{gate} On synthetic data, PS improves over a \emph{df-matched}
LPS where the geometry predicts; $\lambda_{\mathrm{sync}}$ selection is well-founded
via $\mathrm{tr}\,S$ (no grid-boundary railing); $\lambda_{\mathrm{sync}}=0$
reproduces LPS exactly and the path is continuous.\end{gate}
\field{Unblocks} LPS+PS+LDSR; the nerve integration.

%==============================================================================
\section{Composition --- LPS+PS+LDSR}\label{sec:compose}
%==============================================================================

The two axes compose because they are orthogonal: LDSR sets each anchor's chart;
PS couples the anchors' predictions. The natural order is to fix the chart field
first (LDS$\to$LDR/LDSR), then synchronize on it (PS), because PS's overlap graph is
defined by chart overlaps, which the scale field determines. The fully composed
object LPS+PS+LDSR is exactly the cochain picture of the nerve substrate: one
consistency objective over the overlaps combining prediction agreement (PS) and a
coherent chart-geometry field (LDSR). It is built in \S\ref{sec:nerve}, not before,
because it depends on all of its parts.

%==============================================================================
\section{The local covariance base (LCov) and its parallel expansions}\label{sec:lcov}
%==============================================================================

\field{Objective} Build and validate the local conditional covariance / precision /
association on the boundary-safe compositional ($L^\infty$ / Aitchison) geometry ---
the second base estimator --- and carry the \emph{same} refinement algebra onto it.
\field{Inputs} Phase 0; the realistic 16S generators; a boundary-aware compositional
front end (structural-zero and soft-face handling) that this base must establish
first.
\field{Deliverables} Recovery of a known local covariance on synthetic 16S; the
two-layer association object (co-occurrence support $+$ conditional-abundance
precision) with soft-face and $\ell_1$-fusion. Then the parallel expansions:
\textbf{LCov+PS} (synchronize overlapping local covariances --- the second-order
analogue of PS-LPS), \textbf{LCov+LDS/LDR/LDSR} (the \emph{same} chart-geometry
ladder; the dimension/scale machinery is shared with LPS), and \textbf{LCov+PS+LDSR}.
Spec: the chart-aware association design and its addendum.
\field{Gate} \begin{gate} The covariance/precision is recovered on synthetic data
with known structure; the boundary-safe path never reintroduces a pseudocount; the
$\ell_1$ fusion preserves a planted reorganization across a community-state boundary
that $\ell_2$ would smooth away. LCov+LDSR reuses the validated LPS ladder, so its
gate is the recovery gate plus the ladder's reduction/tracking checks.\end{gate}
\field{Unblocks} The nerve integration; the capstone.

\begin{note}
The chart-geometry ladder (Axis B) is genuinely \emph{shared} between LPS and LCov:
both fit something on a local chart whose dimension and scale must be chosen, so the
same $(d_i,K_i)$ field machinery serves both. Only the synchronization analogue (PS
for second-order objects) and the compositional front end are new to LCov.
\end{note}

%==============================================================================
\section{Axis C (cross-cutting) --- grid anchoring}\label{sec:grid}
%==============================================================================

\field{Objective} Decouple the \emph{anchor set} (where local models are centered)
from the \emph{sample}: instead of one local model per data point, center them at a
chosen set of anchors --- a regular or derived grid.
\field{Why it matters} The estimand becomes a function on a controlled anchor set
(resolution chosen, not dictated by sampling density); charts are reusable
(predictable cost); the synchronization graph and the nerve are built on the anchors,
so the anchor choice is upstream of PS and the nerve, not a detail.
\field{Ways to define the grid} A naive ambient lattice is wrong (anchors land in
empty regions); useful constructions all keep anchors \emph{on the data manifold}:
farthest-point / Poisson-disk sampling of the data; a lattice on the estimated
manifold (intrinsic coordinates); coarse-to-fine adaptive refinement driven by the
local-scale field $K_i$ (denser anchors where scale is fine); or the nerve vertices
themselves. The local-scale field of LDSR is the natural driver of an adaptive grid
--- closing a loop between Axis~B and Axis~C.
\field{Gate} \begin{gate} On synthetic data, grid-anchored LPS matches
data-anchored LPS in accuracy at equal effective resolution, with anchors provably
covering the data support (no anchor with an empty or degenerate neighborhood); an
adaptive grid driven by $K_i$ uses fewer anchors at equal accuracy than a uniform
grid.\end{gate}
\field{Unblocks} A cleaner nerve (a controlled cover) and predictable cost for the
capstone on large real data.

%==============================================================================
\section{Cross-cutting --- the synthetic-data and VALENCIA enabler}\label{sec:data}
%==============================================================================

\field{Objective} Supply known-ground-truth, increasingly realistic non-manifold
datasets, because every refinement's gate needs an answer key the real data cannot
give.
\field{Deliverables} Early: the toy stratified generators (the DGP library, now
frozen as Amendment~1) needed by LDS/LDR. Adding for this revision: a
\textbf{varying-scale} generator (a localized region of finer structure / higher
curvature) for the scale ladder, alongside the varying-dimension (G4) and curved
(G3) generators. Maturing: the VALENCIA E1/E2/E3 exploration producing realistic
16S generators for LCov and the capstone. Spec: the implementation notes.
\field{Gate} \begin{gate} Toy and varying-scale generators ready before the ladder's
gates; realistic generators pass a posterior-predictive fidelity check before any
LCov/capstone result built on them is trusted.\end{gate}

\begin{note}
This workstream leads and runs throughout: it is a \emph{dependency} of the ladder,
LCov, and the capstone, so it cannot trail them. The Tier-0 DGP library
(Amendment~1, G1--G7) is the first installment and is already audited and merged.
\end{note}

%==============================================================================
\section{Integration --- the simplicial complex (nerve)}\label{sec:nerve}
%==============================================================================

\field{Objective} Unify, on the nerve of the chart cover, the pieces that the
product structure keeps separate: synchronization (PS), the dimension/scale field
(LDSR), and the covariance (LCov), as one consistency objective over the overlaps.
\field{Inputs} The PS axis, the LDSR rung, the LCov base; the anchoring choice that
defines the cover.
\field{Deliverables} Minimal-first: replace PS's pairwise overlap graph with the
nerve $1$-skeleton of variable-size charts; add $2$-simplices (higher-order
synchronization at junctions); co-regularize chart size and dimension over the nerve
(this \emph{is} LDSR's graph $G$); and, if it pays, the single cochain objective
combining prediction, covariance, and dimension/scale. Spec: the chart-nerve
substrate note.
\field{Gate} \begin{gate} On synthetic data, the nerve $1$-skeleton of variable-size
charts beats the fixed $k$-NN overlap graph on held-out accuracy; $2$-simplex
synchronization either measurably helps at junctions or is dropped; the integrated
objective is validated on known truth.\end{gate}
\field{Unblocks} Phase 3 --- it supplies the basin-and-corridor skeleton.

%==============================================================================
\section{Capstone --- occupancy / quasi-equilibrium / support}\label{sec:capstone}
%==============================================================================

\field{Objective} Find the pure components (recurrent ecological basins) and the
transition phases between them.
\field{Inputs} LCov, the nerve / basin-corridor skeleton, the realistic 16S
generators, and the COT occupation-measure machinery.
\field{Deliverables} Level-set basins on the chart cover; the geometric corridor
discriminator (overlap mass $+$ elongated covariance); the local-covariance basin
shape; alignment of per-subject basins onto shared components; and the transfer
operator on the coarse basin-and-corridor skeleton that earns the
``quasi-equilibrium'' label. Spec: the COT bridge note.
\field{Gate} \begin{gate} On synthetic 16S with \emph{known} components:
geometry-first level-set basins are more stable across scales/bootstraps than the
dictionary-factorization atoms; the transfer-operator metastable count agrees with
the cluster-tree count; the corridor discriminator removes spurious components. Then,
and only then, the HMP/152-subject experiment.\end{gate}
\field{Unblocks} The terminal application.

%==============================================================================
\section{Terminal --- the real 16S / omics application}\label{sec:terminal}
%==============================================================================

\field{Objective} Estimate conditional expectations of binary and continuous
features, and the community occupancy/component structure, on real 16S and other
omics data, with honest uncertainty.
\field{Deliverables} Real-data conditional-expectation surfaces with bands;
occupancy/component decompositions; reported under nested and grouped (subject-level)
cross-validation so the headline numbers are not selection-optimistic or
leakage-inflated.
\field{Gate} \begin{gate} Held-out predictive validity on real data and biological
interpretability of the components --- judged, as throughout, by prediction, not by
intrinsic-geometry scores.\end{gate}

%==============================================================================
\section{Dependencies and decision gates that can re-route the program}\label{sec:dag}
%==============================================================================

The forced order, stated as dependencies:
\begin{enumerate}[leftmargin=1.6em]
\item \textbf{LPS base (Phase 0)} first --- it emits the $S$/df/LOO/band toolkit the
refinements reuse. \emph{(Met.)}
\item \textbf{The two refinement axes develop in parallel on the base.} PS
(\S\ref{sec:ps}) and the dimension/scale ladder (\S\ref{sec:ladder}) are independent;
within the ladder, \textbf{LDS before LDR/LDSR} (regularize an estimator you have).
\item \textbf{LCov base} after Phase 0; it \emph{reuses} the chart-geometry ladder,
so LCov's ladder rungs follow the LPS ones cheaply.
\item \textbf{Anchoring (Axis C)} can be developed alongside, but the \emph{adaptive}
grid depends on the scale field (LDSR), and the nerve depends on the anchor cover.
\item \textbf{Nerve integration} depends on PS, LDSR, and LCov; \textbf{capstone}
depends on the nerve and LCov; \textbf{terminal} depends on all.
\end{enumerate}

\noindent Gates that prune whole branches:
\begin{itemize}[leftmargin=1.4em]
\item \textbf{After LDS.} If the X-only scale already captures the structure (the
y-aware variant adds little Truth-RMSE), keep X-only and the exact toolkit; do not
pay the selection-aware df/band cost.
\item \textbf{After LDR.} If the dimension field does not improve Truth-RMSE, it is
not promoted; a global chart dimension is kept and LDSR reduces to scale-only.
\item \textbf{After LDSR.} If co-regularized scale does not beat scale-frozen LDR at
matched df, the program stays on LDR (fixed scale).
\item \textbf{After the nerve.} If $2$-simplex synchronization does not help, stay on
the $1$-skeleton.
\item \textbf{After the capstone on synthetic.} If geometry-first basins are not more
stable than the existing factorization, the capstone falls back to the COT
dictionary route.
\end{itemize}

%==============================================================================
\section{Immediate next actions (this revision)}\label{sec:next}
%==============================================================================

\begin{enumerate}[leftmargin=1.6em]
\item \textbf{Land the Tier 0--4 integration.} dgp and t4 (uncertainty/bands) are
merged; e19 (bandwidth/CV) and t2 (binary hygiene + the re-opened E0.6) are in their
final audits and merges. This closes the validated LPS base and its toolkit.
\item \textbf{Decide the ladder sequencing.} Build \textbf{LDS} (the joint
dimension/scale estimator) next, choosing the X-only-default / y-aware-opt-in fork;
then choose whether LDR and LDSR land as one joint field (E1.11 absorbs scale) or
staged (LDR, then LDSR/E1.12). Spec in hand: the per-point-scale sketch.
\item \textbf{In parallel}, keep the synthetic enabler ahead: add the varying-scale
generator so the ladder's gates have known-truth inputs.
\end{enumerate}

Once the base is fully landed, the two refinement axes proceed in parallel, LCov
inherits the ladder, the anchoring choice is exercised, and the nerve integrates
them.

%==============================================================================
\begin{thebibliography}{9}\small
\bibitem{plan} \emph{LPS Correctness and Validation: A Precise Experimental Plan}
  (the LPS base / Tier 0--4 spec),
  \nolinkurl{project_briefs/lps_experimental_plan_2026-06-09}.
\bibitem{audit} \emph{LPS/PS-LPS JASA-style audit memo} (PS-LPS linearity, df,
  $\lambda_{\mathrm{sync}}$ selection),
  \nolinkurl{project_briefs/lps_ps_lps_jasa_style_audit_2026-06-09}.
\bibitem{impl} \emph{LPS Local Dimension Regularization: Implementation Notes}
  (the LDR rung; the synthetic generators),
  \nolinkurl{project_briefs/lps_local_dimension_regularization_implementation_notes_2026-06-10}.
\bibitem{scale} \emph{Per-point / co-regularized chart scale for LPS} (the LDS/LDSR
  sketch; the X-only vs y-aware fork),
  \nolinkurl{project_briefs/lps_per_point_scale_sketch_2026-06-13}.
\bibitem{chartaware} \emph{Chart-Aware Local Association Structure on the
  $L^\infty$ Atlas} (the LCov base),
  \nolinkurl{project_briefs/chart_aware_local_association_design_2026-06-09}.
\bibitem{nerve} \emph{The Local-Chart Simplicial Complex as an Organizing Substrate
  for LPS} (the nerve integration),
  \nolinkurl{project_briefs/lps_chart_nerve_substrate_2026-06-10}.
\bibitem{bridge} \emph{Occupancy Supports as Basins on the Chart Cover} (the
  capstone),
  \nolinkurl{vaginal_community_trajectory_types/docs/occupancy_basins_on_chart_cover_bridge_2026-06-10}.
\bibitem{workflow} \emph{Worker--Auditor Research Project Workflow} (the
  audit-independence governance),
  \nolinkurl{~/.codex/notes/workflows/worker_auditor_workflow}.
\end{thebibliography}

\end{document}
